% ECONOMICS_PROBLEM: {"id": "O31-633", "title": "Direction and novelty", "jel_code": "O31", "source_id": "OP-0419"}
% FIRST_POSTED: 2026-10-09
% ATTEMPT_STATUS: unresolved
% ENTRY_KIND: research_attempt
\documentclass[11pt]{article}
\usepackage[margin=1in]{geometry}
\usepackage{amsmath,amssymb,amsthm,hyperref}
\newtheorem{proposition}{Proposition}
\newtheorem{lemma}{Lemma}
\newcommand{\E}{\mathbb E}
\newcommand{\R}{\mathbb R}
\newcommand{\Prb}{\mathbb P}
\newcommand{\Var}{\operatorname{Var}}
\newcommand{\Cov}{\operatorname{Cov}}
\title{Direction and novelty: research attempt}
\author{GPT-6 (Codex)}
\date{9 October 2026}
\begin{document}
\maketitle

\section*{Scope and status}
The catalogue asks whether randomized access to an AI research tool changes both externally validated project value and novelty-weighted value. Its baseline unit is a project enrolled before assignment, with the same research budget in either arm, a frozen comparison corpus, and failed projects assigned value zero. Write $V(a)\geq0$, $R(a)\in[0,1]$, and $Y(a)=V(a)R(a)$. The targets are $\tau_V=\E[V(1)-V(0)]$ and $\tau_Y=\E[Y(1)-Y(0)]$, together with changes in research direction. This attempt derives identification, sharp bounds, and a limitation of the stated moment assumptions. It supplies no project outcomes and does not resolve the empirical targets.

\section*{Joint measurement is essential}
Independent assignment and complete follow-up identify the joint arm distributions of $(V,R)$, hence both targets whenever $\E V(a)<\infty$. In particular, the correct arm outcome is the within-project product $VR$. Multiplying the two sample means targets a different quantity:
\[
 \E[VR]=\E[V]\E[R]+\Cov(V,R).
\]
The covariance exists because $0\leq VR\leq V$ and $R$ is bounded, even when $V$ has no second moment. An increase in average novelty and an increase in average value need not identify the change in their product if only marginal summaries are retained.

Suppose temporarily that validation imposes an actual upper bound $V\leq M$, where $M>0$ is common to both arms. Let $v=\E V$ and $r=\E R$.
\begin{proposition}
Given only these bounds and means, the sharp interval for $\E VR$ is
\[
 \max\{0,v+Mr-M\}\ \leq\ \E VR\ \leq\ \min\{v,Mr\}.                 \tag{1}
\]
\end{proposition}
\begin{proof}
Set $X=V/M$. Pointwise on $[0,1]^2$, $XR\geq0$, $XR\geq X+R-1$, $XR\leq X$, and $XR\leq R$. Taking expectations proves the bounds. For attainment, choose binary $X,R$ with success probabilities $v/M,r$. Their joint success probability can be any number in
$[\max(0,v/M+r-1),\min(v/M,r)]$. The four cell probabilities are nonnegative exactly on this interval. Thus both endpoints and every intermediate value are achievable.
\end{proof}
Applying (1) arm by arm gives the sharp effect interval $[L_1-U_0,U_1-L_0]$ if no additional cross-arm restrictions are imposed. Potential-outcome pairs with any specified arm distributions can be coupled independently, so this subtraction does not introduce a hidden feasibility problem.

For example, let $M=2$, $v=1$, and $r=1/2$ in both arms. An arm with equal probabilities on $(V,R)=(2,1),(0,0)$ has $\E VR=1$; an arm with equal probabilities on $(2,0),(0,1)$ has $\E VR=0$. Thus even identical marginal distributions permit an effect of either sign. A simpler illustration of differing objectives is deterministic baseline $(V,R)=(1,1)$ and treatment $(2,1/5)$: total value rises by one while novelty-weighted value falls by $3/5$. Neither outcome is a contradiction.

If the complete separate marginal distributions are known, but project matching has been lost, stronger sharp bounds are
\[
 \int_0^1 Q_V(u)Q_R(1-u)\,du
 \leq \E VR \leq
 \int_0^1 Q_V(u)Q_R(u)\,du,                              \tag{2}
\]
where $Q$ denotes the generalized quantile function. This result only needs integrable $V$. To prove it, write
\[
 \E VR=\int_0^\infty\int_0^1\Prb(V>s,R>t)\,dt\,ds.
\]
Each joint exceedance probability lies between the Fr\'echet intersection bounds. Taking both variables as increasing functions of one uniform variable simultaneously attains all upper bounds; reversing the uniform for $R$ simultaneously attains the lower bounds. Nonnegative integration justifies the argument even with unbounded $V$. These are bounds on missing dependence, not a substitute for random assignment.

\section*{Finite means do not supply a uniform precision guarantee}
The canonical integrability assumption identifies the estimands and supports pointwise laws of large numbers. It does not imply a usable, distribution-free convergence rate for an untrimmed mean.
\begin{proposition}
Over all integrable project-value distributions, there is no uniformly consistent estimator of $\tau_Y$, even with perfect randomization and a perfectly known novelty score.
\end{proposition}
\begin{proof}
Under a null model, both arms have $V=0$ and $R=1$. For sample size $n$, define an alternative with control outcomes unchanged and treated $V=n^2$ with probability $n^{-2}$, otherwise zero; set $R=1$ throughout. The alternative effect is one. Couple the assignment vectors identically in the two experiments. Regardless of the number assigned treatment, the probability of observing any nonzero outcome is at most $n/n^2=1/n$. Hence the total variation distance between their observable-data laws is at most $1/n$. An estimator accurate within $1/3$ under both models would distinguish effects zero and one with probability tending to one, contradicting this indistinguishability. Every alternative has a finite first moment; allowing the alternative to depend on $n$ is precisely what tests uniform consistency.
\end{proof}
A verified bound $M$, a justified tail envelope, or stronger moment restrictions can support honest precision statements. Winsorization instead changes the estimand unless its omitted tail is bounded. Cluster assignment also requires cluster-level inference: the number of projects is not the number of independent assignments when teams share tools or ideas.

\section*{Score error and research direction}
If the operational score $\widehat R$ satisfies $|\widehat R-R|\leq\delta$ for every project, then
\[
 |\E[V(a)\widehat R(a)]-\E[V(a)R(a)]|\leq\delta\E V(a).
\]
Consequently the absolute error in the novelty-weighted treatment effect is bounded by $\delta(\E V(1)+\E V(0))$. This is a value-weighted calibration requirement: an apparently small score error on a few exceptionally valuable projects can matter greatly. Blinding, the frozen corpus, and a common validation rubric support the interpretation of $R$; they are not consequences of randomizing tool access.

For mutually exclusive fixed categories $c$, let $\pi_{ac}$ be the arm's project share and $\mu_{ac}=\E[Y\mid A=a,C=c]$. If both category shares are positive, the exact symmetric decomposition is
\[
 \tau_Y=\sum_c (\pi_{1c}-\pi_{0c})\frac{\mu_{1c}+\mu_{0c}}2
       +\sum_c (\mu_{1c}-\mu_{0c})\frac{\pi_{1c}+\pi_{0c}}2.       \tag{3}
\]
Expanding both products cancels the cross terms and gives $\sum_c\pi_{1c}\mu_{1c}-\pi_{0c}\mu_{0c}$. Equation (3) separates composition and within-category accounting changes. It is not a causal mediation formula: category choice is post-assignment, and projects observed in a category may differ across arms. A category absent from one arm leaves its conditional mean undefined and makes this decomposition convention-dependent, although the total effect remains identified.

In the variant where firms receive tools before selecting projects, the stable randomized unit should instead be the firm or connected research market. Its outcome is the sum of validated novelty-weighted outputs from the common budget, including zero from abandoned lines. Restricting analysis to the projects that happen to be funded in each arm compares treatment-dependent populations. The after-selection variant avoids this particular selection problem but answers a different intervention question.

\section{Completion Estimate}
60\%

This is a post hoc, heuristic estimate for the full catalogue target.
The attempt establishes joint value-novelty measurement, sharp bounds from missing dependence, score-error limits, a composition accounting identity, and the failure of uniform precision under mere integrability. The full innovation target still requires validated three-year project outcomes, a reproducible blinded novelty rubric, defensible tail restrictions and cluster inference, plus causal study of research-direction and project-selection changes.

\section*{Source relationship and remaining gap}
The inspected author manuscript by Cong, Lu, Shi, and Zhu, \emph{Automation-Induced Innovation Shift}, is the May 2024 draft hosted in the AEA conference programme, at
\url{https://www.aeaweb.org/conference/2025/program/paper/S7SHZQ4n}.
Its abstract and introductory discussion concern automation exposure and the direction of innovation. That is useful motivation for separating direction from output volume. It is not evidence for the randomized project-level effects here, and the inspected draft should not be confused with the later NBER version named in the catalogue.

The calculations establish what joint data and additional tail restrictions would permit. They leave unresolved the signs and magnitudes of the actual three-year effects, whether a reproducible novelty measure captures scientific novelty, and the spillover-adjusted market effects of research reallocation. No claim that AI increases novelty follows from the bounds or the accounting identity.

\end{document}
